\chapter{The Sixteen Modes Applied}
\label{ch:modes-applied}

\begin{versebox}
Sixteen modes at first, then looking\\
at the directions from direction's vantage ---\\
compress it all with the goad,\\
and with three methods, analyze the discourse.
\end{versebox}

\noindent That verse is the key to this entire section --- a mnemonic for the whole system. We have now defined each of the sixteen modes in the abstract. The question is: do they actually work? Take a single passage and run all sixteen modes through it. If the system is real, every mode should yield a reading. If even one mode comes up empty, the system has a gap.

The Netti takes the test.\footnote{Pali: \textit{``Vuttā, tassā niddeso kuhiṃ daṭṭhabbo? Hārasampāte.''} ``These [sixteen modes] have been stated. Where is their demonstration to be seen? In the Modes Applied (\textit{hārasampāta}).'' The word \textit{sampāta} means ``falling together, convergence'' --- the place where all sixteen modes converge on a single passage.}

The test passage is a verse from the Udāna:

\begin{versebox}
With an unguarded mind,\\
struck down by wrong view,\\
overcome by sloth-and-torpor ---\\
one comes under Māra's sway.
\end{versebox}

Four lines. Four problems. And now, sixteen different ways to read them.\footnote{Pali: \textit{``Arakkhitena cittena, micchādiṭṭhihatena ca; thinamiddhābhibhūtena, vasaṃ mārassa gacchatī''ti.''} (Ud 32). The counterpart verse --- the ``bright side'' --- runs: ``Therefore, for one with a guarded mind, whose range is right intention, with right view as forerunner, knowing arising and passing away --- the monk who overcomes sloth-and-torpor should abandon all bad destinations.'' The Netti applies its sixteen modes to both verses, but the bulk of the analysis uses the first (dark) verse as the worked example.}

\section*{Mode 1: The Teaching Mode Applied}

``With an unguarded mind'' --- what is the Buddha teaching here? Heedlessness. That is Māra's foothold.\footnote{Pali: \textit{``Arakkhitena cittenāti kiṃ desayati? Pamādaṃ --- taṃ maccuno padaṃ.''} ``Unguarded mind'' = teaching about heedlessness (\textit{pamāda}), which is ``Māra's foothold'' (\textit{maccuno padaṃ}).}

``Struck down by wrong view'' --- wrong view means seeing permanence where there is none. That is a distortion. And the distortion has a specific structure: it distorts three things --- perception, thought, and view --- across four grounds: seeing the five aggregates as self. Form as self, or a self possessing form, or form within self, or self within form. The same for feeling, perception, formations, and consciousness. Twenty permutations in all.\footnote{Pali: \textit{``Micchādiṭṭhihataṃ nāma vuccati yadā `anicce nicca'nti passati, so vipallāso. So pana vipallāso kiṃlakkhaṇo? Viparītaggāhalakkhaṇo. So kiṃ vipallāsayati? Tayo dhamme --- saññaṃ cittaṃ diṭṭhiṃ. So kuhiṃ vipallāsayati? Catūsu attabhāvavatthūsu --- rūpaṃ attato samanupassati, rūpavantaṃ vā attānaṃ, attani vā rūpaṃ, rūpasmiṃ vā attānaṃ.''} The fourfold distortion applied to each of the five aggregates yields twenty modes of identity-view (\textit{vīsativatthukā sakkāyadiṭṭhi}).}

And then the Netti maps each aggregate to its own specific distortion: form is the ground for ``beauty in the unattractive.'' Feeling is the ground for ``pleasure in suffering.'' Perception and formations are the ground for ``self in what is not-self.'' Consciousness is the ground for ``permanence in the impermanent.'' Two defilements corrupt the mind: craving and ignorance. Craving distorts in two directions --- ``beautiful'' and ``pleasant.'' Ignorance distorts in two directions --- ``permanent'' and ``self.'' The view-distortion looks backward, at what has already been; the craving-distortion looks forward, at what is yet to come.\footnote{Pali: \textit{``Dve dhammā cittassa saṃkilesā --- taṇhā ca avijjā ca. Taṇhānivutaṃ cittaṃ dvīhi vipallāsehi vipallāsīyati `asubhe subha'nti `dukkhe sukha'nti. Diṭṭhinivutaṃ cittaṃ dvīhi vipallāsehi vipallāsīyati `anicce nicca'nti `anattani attā'nti.''} Craving-shrouded mind → beauty-distortion + pleasure-distortion. View-shrouded mind → permanence-distortion + self-distortion. The temporal distinction: view-distortion regards the past aggregates as self; craving-distortion delights in the future aggregates.}

``Overcome by sloth-and-torpor'' --- sloth is the mind's unwieldiness; torpor is the body's sluggishness. ``Coming under Māra's sway'' --- both the Māra of defilement and the Māra that is the being-aggregate itself.

The Teaching Mode's conclusion: the Buddha is teaching two truths here --- suffering and its origin. He teaches the Dhamma for the full understanding of suffering and the abandoning of its origin. What fully understands and what abandons --- that is the path. The abandoning of craving and ignorance --- that is cessation. Four truths in four lines.\footnote{Pali: \textit{``Imāni bhagavatā dve saccāni desitāni --- dukkhaṃ samudayo ca. Tesaṃ bhagavā pariññāya ca pahānāya ca dhammaṃ deseti. Yena ca parijānāti yena ca pajahati, ayaṃ maggo. Yaṃ taṇhāya avijjāya ca pahānaṃ, ayaṃ nirodho. Imāni cattāri saccāni.''} The Teaching Mode reads four Noble Truths out of a four-line verse that never mentions them explicitly. This is the mode's power: it finds the gratification-danger-escape structure latent in any passage. The closing formula: \textit{``Tenāhāyasmā mahākaccāyano `assādādīnavatā'ti.''}}

\section*{Mode 2: The Investigation Mode Applied}

This is the longest and most elaborate application. The Investigation Mode takes the same verse and launches a detailed analysis of the entire path.

It begins with a surprising observation: craving comes in two kinds --- unwholesome and wholesome. Unwholesome craving leads to continued existence. But wholesome craving --- the craving for liberation, the fierce longing that asks ``When will I realize what the noble ones have realized?'' --- that craving leads to dismantling. Even conceit comes in two kinds: the unwholesome kind that generates suffering, and the wholesome kind that uses pride as fuel for abandoning pride.\footnote{Pali: \textit{``Tattha taṇhā duvidhā kusalāpi akusalāpi. Akusalā saṃsāragāminī, kusalā apacayagāminī pahānataṇhā. Mānopi duvidho. Yaṃ mānaṃ nissāya mānaṃ pajahati, ayaṃ māno kusalo.''} The Investigation Mode's first move: even defilements can be wholesome when they become the fuel for their own transcendence. The ``renunciation-based grief'' (\textit{nekkhammasitaṃ domanassaṃ}) --- the sorrow at not yet having attained liberation --- is classified as wholesome craving. This is psychologically astute: the spiritual aspiration that says ``I must get free'' is technically a form of wanting, but it is the wanting that ends all wanting.}

The Investigation Mode then traces the entire Noble Eightfold Path through the fourth absorption. In the fourth absorption, the mind is pure, bright, unblemished, free of defilements, malleable, wieldy, steady, and imperturbable. There, a person attains eight things: six direct knowledges and two special distinctions. And this concentration is six-factored: a mind free of covetousness toward the whole world, free of ill-will toward all beings, with aroused energy, a tranquil body, a concentrated mind, and established mindfulness.\footnote{Pali: \textit{``Catutthe hi jhāne aṭṭhaṅgasamannāgataṃ cittaṃ bhāvayati parisuddhaṃ pariyodātaṃ anaṅgaṇaṃ vigatūpakkilesaṃ mudu kammaniyaṃ ṭhitaṃ āneñjappattaṃ. So tattha aṭṭhavidhaṃ adhigacchati --- cha abhiññā dve ca visese.''} The eight qualities of the fourth-absorption mind, in progressive order: pure → bright → unblemished → free of defilements → malleable → wieldy → steady → imperturbable. Each quality is the condition for the next. The six direct knowledges include psychic powers, the divine ear, mind-reading, recollection of past lives, the divine eye, and the destruction of the taints.}

The three character-types each find liberation through their own door. The lust-character goes out through the signless liberation-door, training in the higher mind, abandoning the unwholesome root of greed --- washing away lust's stain, shaking off lust's dust, vomiting lust's poison, extinguishing lust's fire, pulling out lust's dart, untangling lust's tangle. The hatred-character goes out through the desireless liberation-door, training in the higher virtue, abandoning the unwholesome root of hatred. The delusion-character goes out through the emptiness liberation-door, training in the higher wisdom, abandoning the unwholesome root of delusion.\footnote{Pali: \textit{``Rāgacarito puggalo animittena vimokkhamukhena niyyāti \ldots lobhaṃ akusalamūlaṃ pajahanto \ldots rāgamalaṃ pavāhento rāgarajaṃ niddhunanto rāgavisaṃ vamento rāgaggiṃ nibbāpento rāgasallaṃ uppāṭento rāgajaṭaṃ vijaṭento.''} The vivid sixfold metaphor for each character-type's purification: washing away the stain (\textit{mala}), shaking off the dust (\textit{raja}), vomiting the poison (\textit{visa}), extinguishing the fire (\textit{aggi}), pulling out the dart (\textit{salla}), untangling the tangle (\textit{jaṭā}). The same six metaphors are applied in parallel for hatred and delusion, with their respective liberation-doors switched. Note the unexpected mapping: the \textit{signless} door corresponds to concentration (\textit{samādhikkhandha}), the \textit{desireless} to virtue (\textit{sīlakkhandha}), and \textit{emptiness} to wisdom (\textit{paññākkhandha}).}

The Investigation Mode then runs through the Ten Powers of the Tathāgata --- the ten types of knowledge that make a Buddha a Buddha. The first power is the knowledge of what is possible and what is impossible: a fully awakened Buddha could not have any truth unrealized; a taint-destroyer could not have any taint unextinguished; someone who follows the teaching could not fail to make progress. But equally: a woman could not be a universal monarch, or Sakka, or Māra, or Brahmā, or a Tathāgata. Two Buddhas could not arise simultaneously in a single world-system. Bad conduct could not yield pleasant results; good conduct could not yield unpleasant results.\footnote{Pali: \textit{``Sammāsambuddhassa te sato ime dhammā anabhisambuddhāti netaṃ ṭhānaṃ vijjati \ldots Itthī tathāgato arahaṃ sammāsambuddho siyāti netaṃ ṭhānaṃ vijjati \ldots''} The ``possible and impossible'' (\textit{ṭhānāṭṭhāna}) analysis, the first of the ten Tathāgata-powers. The passage includes the controversial claim about gender and Buddhahood. \textbf{Translator's note:} This passage reflects the social assumptions of its ancient Indian context. The Netti is quoting standard Abhidhamma doctrine here, not making an original argument.}

The second power is the knowledge of the destinations of all actions: beings who commit the three kinds of misconduct go to the lower realms; beings who commit the three kinds of good conduct go to the happy realms; those who develop the path attain final extinguishment. The third through tenth powers cover: the diversity of elements and realms; the diversity of beings' inclinations; the four types of action and their results; the defilement and purification of absorptions, liberations, concentrations, and attainments; the maturity of beings' faculties; the recollection of past lives; the divine eye that sees beings passing away and arising; and the destruction of all taints.\footnote{Pali: \textit{``Iti ṭhānāṭṭhānagatā sabbe khayadhammā \ldots Sabbaññutā pattā \ldots idaṃ bhagavato dasamaṃ balaṃ sabbāsavaparikkhayaṃ ñāṇaṃ.''} The ten powers of the Tathāgata (\textit{dasabalañāṇa}) are given in abbreviated form, with the first and second developed at length. Each power is a specific kind of knowledge that enables the Buddha to teach in a way perfectly calibrated to each being's condition. The Netti's point is that the Investigation Mode, applied to a single verse, can expand to encompass the entire scope of the Buddha's teaching capacity.}

The verse ``All beings will die'' receives the same treatment: the Investigation Mode reads it as teaching three rounds (the round of suffering, the round of action, and the round of defilement), three types of defilement (craving, view, misconduct), three types of purification (calm, insight, good conduct), and three paths (the demeritorious, the meritorious, and the path that transcends both merit and demerit).\footnote{Pali: \textit{``Sabbe sattā marissanti, maraṇantaṃ hi jīvitaṃ \ldots''} (SN 1.133). The Investigation Mode's reading: ``All beings will die \ldots to hell through evil deeds'' = defilement side (three rounds: suffering, action, defilement). ``Good deeds to a happy destination'' = merit side. ``Others, having developed the path, attain final extinguishment without taints'' = the transcending of both merit and demerit. The three types of defilement are then mapped to three purifications: craving-defilement → calm (\textit{samatha}), view-defilement → insight (\textit{vipassanā}), misconduct-defilement → good conduct (\textit{sucarita}).}

\section*{Modes 3--16: The Rapid Applications}

The remaining fourteen modes are applied with startling economy. Each takes the same counterpart verse --- ``Therefore, for one with a guarded mind, whose range is right intention, with right view as forerunner, knowing arising and passing away, the monk who overcomes sloth-and-torpor should abandon all bad destinations'' --- and extracts its own distinctive reading in a few sentences.\footnote{Pali: \textit{``Tasmā rakkhitacittassa, sammāsaṅkappagocaro; sammādiṭṭhipurekkhāro, ñatvāna udayabbayaṃ; thinamiddhābhibhū bhikkhu, sabbā duggatiyo jahe''ti.''} This is the ``bright'' counterpart verse, reversing each element of the Māra verse: unguarded → guarded; wrong view → right view; sloth-and-torpor overcome; coming under Māra's sway → abandoning all bad destinations. Modes 3--16 all use this verse as their test case.}

The \textbf{Fitness Mode} checks whether each phrase logically follows from the last. A guarded mind will have right intention as its range --- that fits. Right intention will lead to right view --- that fits. Right view as forerunner will lead to knowing arising and passing away --- that fits. Knowing arising and passing away will lead to abandoning all bad destinations --- that fits. Every link holds.\footnote{Pali: \textit{``Rakkhitacittassa sammāsaṅkappagocaro bhavissatīti yujjati, sammāsaṅkappagocaro sammādiṭṭhi bhavissatīti yujjati \ldots sabbā duggatiyo jahanto sabbāni duggativinipātabhayāni samatikkamissatīti yujjatīti.''} The Fitness Mode applied: five ``it fits'' (\textit{yujjati}) statements, one per phrase-transition. The chain is logically coherent.}

The \textbf{Basis Mode} traces the foundation of each phrase: ``guarded mind'' is the basis of the three kinds of good conduct. ``Right intention'' is the basis of calm. ``Right view as forerunner'' is the basis of insight. ``Knowing arising and passing away'' is the basis of the vision-ground. ``Overcoming sloth-and-torpor'' is the basis of energy. ``Abandoning all bad destinations'' is the basis of development.\footnote{Pali: \textit{``Tasmā rakkhitacittassāti tiṇṇaṃ sucaritānaṃ padaṭṭhānaṃ. Sammāsaṅkappagocaroti samathassa padaṭṭhānaṃ. Sammādiṭṭhipurekkhāroti vipassanāya padaṭṭhānaṃ. Ñatvāna udayabbayanti dassanabhūmiyā padaṭṭhānaṃ. Thinamiddhābhibhū bhikkhūti vīriyassa padaṭṭhānaṃ. Sabbā duggatiyo jaheti bhāvanāya padaṭṭhānaṃ.''}}

The \textbf{Characteristic Mode} shows that when right view is grasped, the entire Eightfold Path is grasped --- because right view gives rise to right intention, which gives rise to right speech, right action, right livelihood, right effort, right mindfulness, right concentration, right liberation, and the knowledge-and-vision of liberation. Ten links, each flowing from the one before.\footnote{Pali: \textit{``Sammādiṭṭhiyā gahitāya gahito bhavati ariyo aṭṭhaṅgiko maggo. Taṃ kissa hetu? Sammādiṭṭhito hi sammāsaṅkappo pabhavati \ldots sammāvimuttito sammāvimuttiñāṇadassanaṃ pabhavati.''} The Characteristic Mode shows that grasping one factor (right view) entails grasping all ten factors of the tenfold path (the eightfold path plus right liberation and right knowledge-and-vision). This is the ``one grasped, all grasped'' principle.}

The \textbf{Fourfold Array} asks what the Buddha's actual intention is: those who wish to be freed from bad destinations should live by the Dhamma. The story of Kokālika is the counter-example --- he corrupted his mind toward the elders Sāriputta and Moggallāna and was reborn in the Great Lotus Hell.\footnote{Pali: \textit{``Kokāliko hi sāriputtamoggallānesu theresu cittaṃ padosayitvā mahāpadumaniraye upapanno.''} The Fourfold Array applied: the Buddha's actual intention behind the verse is a warning --- guard your mind or face the consequences. Kokālika's story is the vivid illustration.}

The \textbf{Conversion Mode} reads the verse as the Four Noble Truths: ``guarded mind'' and ``right intention'' = calm. ``Right view as forerunner'' = insight. ``Knowing arising and passing away'' = comprehension of suffering. ``Overcoming sloth-and-torpor'' = abandoning the origin. ``Abandoning all bad destinations'' = cessation.\footnote{Pali: \textit{``Tasmā rakkhitacittassa sammāsaṅkappagocaroti samatho. Sammādiṭṭhipurekkhāroti vipassanā. Ñatvāna udayabbayanti dukkhapariññā. Thinamiddhābhibhū bhikkhūti samudayapahānaṃ. Sabbā duggatiyo jaheti nirodho. Imāni cattāri saccāni.''}}

The \textbf{Analysis Mode}: the wholesome side should be analyzed as wholesome; the unwholesome side as unwholesome. The \textbf{Reversal Mode}: calm and insight, once developed, yield cessation as their fruit; suffering is comprehended; the origin is abandoned; the path is developed as counterpart. The \textbf{Synonym Mode} unpacks each term into its synonyms: ``mind'' = mind, mental organ, consciousness, mind-faculty, mind-base; ``right intention'' = renunciation-intention, non-ill-will intention, non-cruelty intention; ``right view'' = wisdom-sword, wisdom-blade, wisdom-jewel, wisdom-lamp, wisdom-goad, wisdom-palace.\footnote{Pali: \textit{``Cittaṃ mano viññāṇaṃ manindriyaṃ manāyatanaṃ \ldots Sammādiṭṭhi nāma paññāsatthaṃ paññākhaggo paññāratanaṃ paññāpajjoto paññāpatodo paññāpāsādo.''} The Synonym Mode's reading: six synonyms for ``mind,'' three for ``right intention,'' six metaphorical synonyms for ``right view.'' Wisdom is simultaneously a weapon, a blade, a jewel, a light, a goad, and a palace.}

The \textbf{Designation Mode} labels each phrase by its function: ``guarded mind'' designates the basis of mindfulness; ``right intention'' designates the development of calm; ``right view knowing arising and passing away'' designates the placing of the vision-ground; ``overcoming sloth-and-torpor'' designates the complete abandoning of the origin; ``abandoning all bad destinations'' designates the development of the path.\footnote{Pali: \textit{``Tasmā rakkhitacittassāti padaṭṭhānapaññatti satiyā. Sammāsaṅkappagocaroti bhāvanāpaññatti samathassa. Sammādiṭṭhipurekkhāro ñatvāna udayabbayanti dassanabhūmiyā nikkhepapaññatti \ldots''}}

The \textbf{Entry Mode} enters through all five doorways: right view grasps the five faculties (entry by faculties); those faculties are true knowledge, and with the arising of true knowledge comes the cessation of ignorance, and with the cessation of ignorance the cessation of formations, and so on (entry by dependent origination); those five faculties are gathered into three aggregates (entry by aggregates); those aggregates belong to the dhamma-element (entry by elements); and the dhamma-element belongs to the dhamma-base (entry by sense-bases).\footnote{Pali: \textit{``Sammādiṭṭhiyā gahitāya gahitāni bhavanti pañcindriyāni, ayaṃ indriyehi otaraṇā. Tāniyeva indriyāni vijjā, vijjuppādā avijjānirodho \ldots ayaṃ paṭiccasamuppādena otaraṇā. Tāniyeva pañcindriyāni tīhi khandhehi saṅgahitāni \ldots ayaṃ khandhehi otaraṇā.''} Five doorways, all from a single verse. The Entry Mode compresses the entire analytical framework into a cascade: faculties → dependent origination → aggregates → elements → sense-bases.}

The \textbf{Clearing Mode}: where the beginning is clean, the question is answered; where the beginning is not clean, the question is not yet answered. The \textbf{Resolution Mode} distinguishes ``oneness'' from ``difference'' within each term: ``mind'' is the general category (oneness); mind, mental organ, consciousness, mind-faculty, mind-base are the specific distinctions (difference). ``Right intention'' is the oneness; renunciation, non-ill-will, and non-cruelty are the difference. ``Arising and passing away'' is the oneness; the twelve links of dependent origination running forward and backward are the difference. ``Bad destinations'' is the oneness; bad destinations relative to gods and humans are one thing, but relative to nirvana \textit{all} destinations are bad --- and that is the deeper difference.\footnote{Pali: \textit{``Sabbā duggatiyo jaheti ekattatā. Devamanusse vā upanidhāya apāyā duggati, nibbānaṃ vā upanidhāya sabbā upapattiyo duggati, ayaṃ vemattatā.''} The Resolution Mode's most striking move: ``bad destinations'' has two levels. Relative to the heavens, only the lower realms are bad. Relative to nirvana, \textit{every} rebirth is a bad destination. The deeper reading swallows the shallower one.}

The \textbf{Requisites Mode}: the entire verse is the requisite of calm and insight. The \textbf{Integration Mode} closes the demonstration by tracing the chain of integration one final time: a guarded mind leads to the three kinds of good conduct; right view, once developed, develops the entire Noble Eightfold Path; the Eightfold Path leads to right liberation; right liberation leads to the knowledge-and-vision of liberation. This is the person without remainder. This is the nirvana-element without remainder.\footnote{Pali: \textit{``Sammādiṭṭhiyā bhāvitāya bhāvito bhavati ariyo aṭṭhaṅgiko maggo \ldots sammāvimuttito sammāvimuttiñāṇadassanaṃ pabhavati. Ayaṃ anupādiseso puggalo anupādisesā ca nibbānadhātu.''} The Integration Mode's final reading: the verse, fully integrated, points to the person without remainder (\textit{anupādiseso puggalo}) and the nirvana-element without remainder (\textit{anupādisesā nibbānadhātu}). This is the ultimate destination of the entire hermeneutical system --- liberation.}

\ornament

Sixteen modes. One verse. Every mode yields a reading. The system has no gaps.

That is what the Venerable Mahākaccāyana meant: ``Sixteen modes at first, then looking at the directions from direction's vantage --- compress it all with the goad, and with three methods, analyze the discourse.''\footnote{Pali: \textit{``Soḷasa hārā paṭhamaṃ, disalocanato disā viloketvā; saṅkhipiya aṅkusena hi, nayehi tīhi niddise suttanti.''} The verse that opened this section now closes it, with its meaning demonstrated rather than merely asserted. The ``three methods'' (\textit{nayehi tīhi}) will be the subject of the next chapter.}
